\section{周期八的精确谱边}\label{sec:period-eight-edge}

第~\ref{sec:periodic-construction} 节中的有理数便于证明严格有限反例，但它并非
真实谱边。本节证明定理~\ref{thm:exact-period-eight-edge}。

\begin{figure}[H]
\centering
\begin{tikzpicture}[
  scale=1.05,
  defect/.style={circle,draw,fill=black,minimum size=6.5mm,inner sep=0pt},
  regular/.style={circle,draw,fill=white,minimum size=6.5mm,inner sep=0pt},
  signplus/.style={font=\small,text=white},
  signminus/.style={font=\small,text=black}
]
  \foreach \i/\ang/\sgn in {
    0/90/+,
    1/45/-,
    2/0/-,
    3/-45/-,
    4/-90/+,
    5/-135/-,
    6/180/-,
    7/135/-
  }{
    \coordinate (p\i) at (\ang:1.65);
    \ifnum\i=0
      \node[defect] (n\i) at (p\i) {};
      \node[signplus] at (p\i) {\(\sgn\)};
    \else\ifnum\i=4
      \node[defect] (n\i) at (p\i) {};
      \node[signplus] at (p\i) {\(\sgn\)};
    \else
      \node[regular] (n\i) at (p\i) {};
      \node[signminus] at (p\i) {\(\sgn\)};
    \fi\fi
    \node[font=\scriptsize] at (\ang:2.12) {\(Q_{\i}\)};
  }
  \draw[thin] (n0)--(n1)--(n2)--(n3)--(n4)--(n5)--(n6)--(n7)--cycle;
  \draw[{Latex[length=2mm]}-{Latex[length=2mm]},line width=0.75pt]
    ($(p0)+(0.20,0)$)--node[left=10pt,font=\small]{相距 \(4\)} ($(p4)+(0.20,0)$);
  \node[anchor=west,align=left,font=\small] at (2.55,0.45)
    {\(\tau=(+,+,-,+,-,-,+,-)\)\\[3pt]
     \(Q=(+,-,-,-,+,-,-,-)\)\\[3pt]
     \(D(Q)=\{0,4\}\)};
\end{tikzpicture}

\caption{目标周期八相位。三角形通量字 \(\tau\) 诱导四边形通量字 \(Q\)；
其中两个正号以黑色实心点标出，并在八点胞中互为对径，故在平移与反射意义下
\(D(Q)=\{0,4\}\)。第~\ref{sec:eight-barrier} 节的三分定理证明这正是低于阈值
8 的情形。}
\label{fig:period8-target}
\end{figure}

\subsection{端点根}

当 \(c=2\) 时，式~\eqref{eq:floquet-polynomial} 化为

\begin{equation}
\label{eq:edge-polynomial}
P(y,2)=y^{4}-16y^{3}+76y^{2}-96y+16.
\end{equation}

作平移 \(y=x+4\)，则

\begin{equation}
\label{eq:edge-depressed-quartic}
P(x+4,2)=x^{4}-20x^{2}+80.
\end{equation}

令 \(w=x^2\)，得到 \(w^2-20w+80=0\)，其根为
\(10\pm2\sqrt5\)。式~\eqref{eq:edge-polynomial} 的四个实根按递增次序为

\begin{equation}
\label{eq:edge-four-roots}
\begin{aligned}
&4-\sqrt{10+2\sqrt{5}}, \\
&4-\sqrt{10-2\sqrt{5}}, \\
&4+\sqrt{10-2\sqrt{5}}, \\
&4+\sqrt{10+2\sqrt{5}}
\end{aligned}
\text{。}
\end{equation}

因此最大的端点根为

\begin{equation}
\label{eq:eta-definition}
\eta=4+\sqrt{10+2\sqrt{5}}.
\end{equation}

\subsection{精确上界与等号唯一性}

令

\begin{equation}
\begin{aligned}
s&=\sqrt{10+2\sqrt{5}},       \eta=4+s, \\
u&=y-\eta,                   t=2-c
\end{aligned}
\text{，}
\end{equation}

代入式~\eqref{eq:floquet-polynomial} 得恒等式

\begin{equation}
\label{eq:eta-positive-expansion}
\begin{aligned}
&P(\eta+u,2-t) \\
&=u^{4}+4s u^{3}+2u^{2}t+(40+12\sqrt{5})u^{2} \\
&\quad +4s ut+8\sqrt{5}s u+t^{2}+(4\sqrt{5}-3)t
\end{aligned}
\text{。}
\end{equation}

所有系数均为正。因此，当 \(u,t\geq0\) 时，右端非负，且仅在 \(u=t=0\)
时为零。由于每个单位圆纤维都有 \(c\leq2\)，式~
\eqref{eq:eta-positive-expansion} 排除大于 \(\eta\) 的每个平方特征值，并且只
在 \(c=2\) 时允许等号。

当 \(c=2\) 时，\(\eta\) 是式~\eqref{eq:edge-polynomial} 的根，所以
\(\pm\sqrt\eta\) 是 \(H(1)\) 的特征值。因此

\begin{equation}
\label{eq:exact-infinite-edge}
\sup_{\lvert z\rvert=1} \rho(H(z))^{2}=\eta.
\end{equation}

在单位圆上，\(z+z^{-1}=2\) 当且仅当 \(z=1\)。这同时证明了定理~
\ref{thm:exact-period-eight-edge} 的数值与唯一性断言。

为后文使用，还记录 \(\eta\) 的最小多项式为

\begin{equation}
\label{eq:eta-minimal-polynomial}
Y^{4}-16Y^{3}+76Y^{2}-96Y+16.
\end{equation}

事实上，平移 \(Y=X+4\) 把式~\eqref{eq:eta-minimal-polynomial} 化为
\(X^4-20X^2+80\)。由 Gauss 引理，有理因式分解可取为首一整系数二次式。
三次项和一次项消失，使它只能是 \((X^2+a)(X^2+b)\)，其中
\(a+b=-20\)、\(ab=80\)，或 \((X^2+rX+s)(X^2-rX+s)\)，其中
\(s^2=80\)。第一种情形会使 \(a,b\) 成为 \(T^2+20T+80\) 的根，但其判别式
80 不是平方数；第二种情形不存在有理数 \(s\)。故式~
\eqref{eq:eta-minimal-polynomial} 在 \(\mathbb Q\) 上不可约。数 \(\eta\) 位于
\((\frac{1951}{250},\frac{1561}{200})\)；特别地，\(\eta<8\)。

\subsection{最高谱带的单调性}

令 \(r(c)\) 为 \(P(y,c)\) 的最大实根。当 \(c=-2\) 时，

\begin{equation}
\label{eq:minus-one-fiber-factorization}
\begin{aligned}
P(y,-2)&=(y^{2}-12y+34)(y^{2}-4y+2), \\
r(-2)&=6+\sqrt{2}=:y_{0}
\end{aligned}
\text{。}
\end{equation}

当 \(-2<c\leq2\) 时，代入得

\begin{equation}
\label{eq:minus-one-fiber-comparison}
P(y_{0},c)=(c+2)(c+5-8\sqrt{2})<0.
\end{equation}

由于 \(P(y,c)\) 随 \(y\) 趋于正无穷，式~
\eqref{eq:minus-one-fiber-comparison} 蕴含当 \(c>-2\) 时
\(r(c)>y_0\)。此外，

\begin{equation}
\label{eq:floquet-c-derivative}
P_c(y,c)=2(c-c_{0}(y)),       c_{0}(y)=y^{2}-8y+\frac{13}{2}.
\end{equation}

函数 \(c_0\) 在 \(y\geq y_0>4\) 上递增，且

\begin{equation}
c_{0}(y_{0})=-\frac{7}{2}+4\sqrt{2}>2.
\end{equation}

因此，在 \(y\geq y_0\)、\(c\leq2\) 的整个区域内都有 \(P_c(y,c)<0\)。若
\(-2\leq c_1<c_2\leq2\) 且 \(y_1=r(c_1)\)，则
\(P(y_1,c_2)<P(y_1,c_1)=0\)，所以 \(P(\,\cdot\,,c_2)\) 在 \(y_1\) 之上
有根。故 \(r(c)\) 在 \([-2,2]\) 上严格递增。

\subsection{有限图的和乐扇区}

当 \(\alpha=+1\) 时，有限集合 \(z^L=1\) 包含 \(z=1\)。由式~
\eqref{eq:exact-infinite-edge}，

\begin{equation}
\label{eq:positive-holonomy-edge}
\rho(A_{8L,+1})^{2}=\eta.
\end{equation}

当 \(\alpha=-1\) 时，允许参数为
\(z_k=\exp((2k+1)\pi i/L)\)。其中最大的 \(c\) 值是
\(2\cos(\frac{\pi}{L})\)，严格小于二。刚证明的单调性给出

\begin{equation}
\label{eq:negative-holonomy-edge}
\rho(A_{8L,-1})^{2}=r(2\cos(\frac{\pi}{L}))<\eta.
\end{equation}

当 \(L\) 趨于无穷时，\(2\cos(\frac{\pi}{L})\) 趨于二；Hermitian 纤维的
特征值连续性使式~\eqref{eq:negative-holonomy-edge} 收敛到 \(\eta\)。所以两种
和乐扇区具有相同的无限体积谱边，但每个有限尺度上只有正扇区达到该值。
